% Generated draft: requires mathematical and source review.
\documentclass[11pt]{article}
\usepackage[margin=25mm]{geometry}
\usepackage{fontspec}
\setmainfont{Latin Modern Roman}
\usepackage{amsmath,amssymb,mathtools,unicode-math}
\setmathfont{Latin Modern Math}
\usepackage{enumitem,needspace,xurl,hyperref}
\setlength{\parindent}{0pt}
\setlength{\parskip}{5pt}
\setlength{\emergencystretch}{3em}
\newcommand{\N}{\mathbb{N}}
\newcommand{\Z}{\mathbb{Z}}
\newcommand{\Q}{\mathbb{Q}}
\newcommand{\R}{\mathbb{R}}
\newcommand{\C}{\mathbb{C}}
\newcommand{\F}{\mathbb{F}}
\newcommand{\A}{\mathbb{A}}
\newcommand{\PP}{\mathbb{P}}
\newcommand{\E}{\mathbb{E}}
\newcommand{\eps}{\varepsilon}
\newcommand{\dd}{\,\mathrm{d}}
\begin{document}
\section*{UnsolvedMath ID: 30006309}
\textbf{Source problem code:} catalogue entry
\subsection{Definitions and mathematical statement}
For a specified split algebraic torus \(T\) over a field \(k\), let \(M=\operatorname{Hom}(T,\mathbb G_m)\) be its character lattice. If a nonzero polynomial \(F\) belongs to a specified finite-dimensional \(T\)-representation and has weight decomposition
\[
F=\sum_{\chi\in M}F_\chi,\qquad t\cdot F_\chi=\chi(t)F_\chi,
\]
its weight polytope is
\[
W_T(F)=\operatorname{conv}\{\chi\in M:F_\chi\ne0\}\subset M\otimes_{\Z}\R.
\]
This definition requires the representation, including the action on polynomial coordinates, to be fixed.

The supplied problem asks for a purely combinatorial proof of the coincidence
\[
W\!\left(\Delta_{X\times\mathbb P^{n-1}}\right)
=
W\!\left(\operatorname{Hu}_A\right),
\]
where the source designates \(\Delta_{X\times\mathbb P^{n-1}}\) as a discriminant and \(\operatorname{Hu}_A\) as a Hurwitz form. The requested object is a combinatorial proof of this equality, rather than a new assertion that the equality holds.

\emph{Unresolved scope.} The supplied record does not define \(A\), its relation to \(X\), the base field, the dimension and embedding of \(X\), the meaning or range of \(n\), the embedding of \(X\times\mathbb P^{n-1}\), or the hypotheses ensuring that the designated discriminant and Hurwitz form exist as nonzero defining polynomials. The projective-space notation requires \(n\) to be a positive integer, but the intended range is not supplied. The torus actions and the identification of the weight spaces in which the displayed equality is asserted are also missing. Accordingly, the equation records the intended target; it does not specify an unconditional theorem for all varieties or all configurations. A fully quantified formulation requires recovery of these data from the source context. No definition of ``purely combinatorial proof'' or permitted proof ingredients is supplied.
\subsection{Short English statement}
The problem requests a purely combinatorial proof that the discriminant of \(X\times\mathbb P^{n-1}\) and the Hurwitz form associated with \(A\) have the same weight polytope. The supplied curation says that the equality is established and labels the requested combinatorial reproof open. The geometric hypotheses and the common weight-space identification are missing, so the precise scope requires human review.
\subsection{Sources}
\begin{itemize}
\item Local UnsolvedMath registry, supplied complete source record, permanent ID \(30006309\), source problem code \texttt{catalogue entry}; collection: Oberwolfach Reports Open Problems. The record gives proposed year 2025 and supplies no proposer attribution.
\item \emph{Toric Geometry} (2025), Oberwolfach Reports, pp. 883--940, DOI \texttt{10.4171/owr/2025/19}. Supplied source and literature link: \url{https://doi.org/10.4171/owr/2025/19}.
\item Sano, \emph{Weight Polytopes and Energy Functionals of Toric Varieties}. Supplied literature link: \url{https://doi.org/10.1007/s42543-023-00079-z}.
\item Supplied literature assessment, checked on \texttt{2026-08-22}: the equality of the Hurwitz-form and discriminant weight polytopes is described as established, including by analytic and secondary-polytope methods, while a purely combinatorial reproof is described as explicitly left as a problem. The supplied status is \texttt{open}. This is supplied dated curation derived from the OWR record, not independently verified source reading or a new literature assessment; its time-sensitive status and cited claims require independent verification.
\item The original extraction contains an isolated ambient-projective-space fragment rendered as \(\mathbb P^{(N+1)n-1}\), and says that one can compute \(W(\Delta_{X\times\mathbb P^{n-1}})\) instead of \(W(\operatorname{Hu}_A)\), following reference [3]. The surrounding embedding context and the identity of [3] are not supplied. The extraction also acknowledges support by JSPS KAKENHI Grant Number \texttt{22K03325}; this does not identify the proposer.
\item The registry labels the statement \texttt{corrected\_verified}, with the qualification that only clear formula delimiters, notation, grammar, and typography were repaired and that mathematical content and scope were preserved. This supplied editorial label does not establish mathematical completeness or independent verification.
\end{itemize}
\end{document}
